\documentclass[10pt,a4paper]{amsart}
\usepackage[margin=19mm]{geometry}
\usepackage{amsmath,amssymb,mathtools}
\usepackage[scr=rsfs,cal=euler]{mathalfa}
\usepackage{xfrac}
\usepackage[unicode,hidelinks,bookmarks=false]{hyperref}
\newcommand{\ie}{i.e.\ }
\newcommand{\cf}{cf.\ }
\newcommand{\resp}{respectively\ }
\newcommand{\Subsection}{\subsection}
\providecommand{\nobreakdash}{\nobreak\hbox{-}}
\setlength{\emergencystretch}{2em}
\multlinegap0pt
\usepackage{bm}
\usepackage{bbm}
\usepackage{dsfont}
\usepackage{xspace}
\usepackage{cancel}
\usepackage{mathtools}
\usepackage{amsthm}
\makeatletter
\@ifundefined{Definition}{\newtheorem{Definition}{Definition}}{}
\@ifundefined{Lemma}{\newtheorem{Lemma}[Definition]{Lemma}}{}
\makeatother
\title[On the sharpness of the $C^1$-norm threshold for perturbatio]{On the sharpness of the $C^1$-norm threshold for perturbations in the normally hyperbolic invariant manifold theorem---a toy model perspective}
\author{Yujie Huang, Junhao Li, Lin Wang}
\date{Source: arXiv:2608.04862v1 (CC BY 4.0)}
\begin{document}
\maketitle

\noindent\emph{Source and licence.} On the sharpness of the $C^1$-norm threshold for perturbations in the normally hyperbolic invariant manifold theorem---a toy model perspective. Version arXiv:2608.04862v1 (CC BY 4.0), distributed under the Creative Commons Attribution 4.0 International licence, as recorded in the arXiv licence metadata for this exact version and shown on the abstract page https://arxiv.org/abs/2608.04862v1. Full rights notice: licenses/arxiv-2608.04862-CC-BY-4.0.txt.\par\smallskip

\noindent\textit{Editorial scope.} Counted unit: the Lemma whose source label is \texttt{lem:nh-iff} in the article (source0.tex lines 689--698, source label \texttt{lem:nh-iff}), delivered together with the article's own complete original proof of it (source0.tex lines 700--872, one \texttt{proof} environment of 173 lines). No mathematics is written, repaired or re-derived: statement and proof are the article's own text line for line. Required context retained uncounted: none. Every definition, display and label of those lines is retained unchanged. Omitted from this package and not redistributed: 1--688, 699--699, 873--1844. Nothing inside a retained excerpt is omitted or altered: every retained body is delivered exactly as the article's lines are written, so no omission receipt of any kind accompanies this package. Citations: the retained excerpts carry no citation command at all, so no bibliography entry of the article is transcribed and no reference is redistributed. Reference adaptation: none was needed and none was made; every reference inside the delivered unit resolves inside it, so no unresolved reference or citation is produced. Printed slips kept literal: no printed word, symbol, hypothesis or number of the counted statement or of its proof was altered. Layout and class: the article's own class is \texttt{amsart} with packages including inputenc, amsfonts, amssymb, amsmath, amscd, euscript, verbatim, array, geometry, anysize, fancyhdr, indentfirst, graphicx, color; the wrapper is the lane's pinned \texttt{amsart} editorial wrapper at 10pt on A4, which restates the theorem-like environments the retained text needs and adds the article's own macro definitions that the retained text needs (none), with extra wrapper packages (bm, bbm, dsfont, xspace, cancel, mathtools). The delivered numbering is the unit's own. Assets: no retained span contains a figure environment, a graphics-inclusion command, a PDF-page inclusion, a table or a TikZ picture; no image file is copied into this package, the package declares no asset dependency and no image file is redistributed with it. Licence: version arXiv:2608.04862v1 of this article is distributed under the Creative Commons Attribution 4.0 International licence, as recorded in the arXiv licence metadata for this exact version and shown on the abstract page https://arxiv.org/abs/2608.04862v1. A full rights notice is delivered at licenses/arxiv-2608.04862-CC-BY-4.0.txt. Characters: every retained excerpt is pure ASCII, so the delivered engine, XeLaTeX with its default Latin Modern text font, reports no missing character.
	\begin{Lemma}\label{lem:nh-iff}Fix $r\geq 1$.
		Let $\lambda\in(0,1)$ and assume that
		$F$
		admits a $C^r$ invariant graph $\Gamma=\{(x,\Psi(x))\mid x\in\mathbb{T}\}$.
		Then $\Gamma$ is $r$-normally hyperbolic for $F$ if and only if
		\[
		g'(x)>\lambda^{\frac{1}{r+1}}\qquad\forall x\in\mathbb{T}.
		\]
		where $g(x)=x+\alpha_1+\lambda\Psi(x)+\phi(x)$.
	\end{Lemma}

	\begin{proof}
		We prove the two implications separately.
		
		\noindent\textbf{Necessity ($\Rightarrow$).}
		Assume $\Gamma$ is $r$-normally hyperbolic. Differentiating the invariance equation $F(x,\Psi(x))=(g(x),\Psi(g(x)))$ with respect to $x$ gives
\[
DF(x,\Psi(x))\begin{pmatrix} 1 \\ \Psi'(x) \end{pmatrix}
= \begin{pmatrix} g'(x) \\ \Psi'(g(x))g'(x) \end{pmatrix}
= g'(x)\begin{pmatrix} 1 \\ \Psi'(g(x)) \end{pmatrix}.
\]
Thus the tangent vector field $e_1(x)=(1,\Psi'(x))^{\mathsf{T}}$ satisfies
$DF e_1 = g'\, e_1\circ g$, and the tangential multiplier is $\tau(x)=g'(x)>0$. It remains to compute the stable multiplier.

\medskip

For each $x\in\mathbb{T}$ consider the tangent space $T_{(x,\Psi(x))}(\mathbb{T}\times\mathbb{R}) \simeq \mathbb{R}^2$
and its subspace $T_x\Gamma = \operatorname{span}\{(1,\Psi'(x))\}$.  The normal space
$N_x = T_{(x,\Psi(x))}(\mathbb{T}\times\mathbb{R}) / T_x\Gamma$ is one-dimensional.
Define the linear map $\pi_x: \mathbb{R}^2 \to \mathbb{R}$ by
\[
\pi_x(u,v) = v - u\,\Psi'(x).
\]
Clearly $\ker\pi_x = \{(u,v): v=u\Psi'(x)\} = T_x\Gamma$, and $\pi_x$ is surjective.  It follows from the first isomorphism theorem  that $\pi_x$ induces a natural
linear isomorphism $N_x\simeq\mathbb{R}$, still denoted by $\pi_x$,
which sends the equivalence class $[u,v]$ to $v-u\Psi'(x)$.

The derivative $DF(x,\Psi(x))$ maps $T_x\Gamma$ into $T_{g(x)}\Gamma$,
hence it induces a linear map $\overline{DF}(x): N_x \to N_{g(x)}$ characterized by
\[
\overline{DF}(x) \circ \pi_x = \pi_{g(x)} \circ DF(x).
\]
We now compute $\overline{DF}(x)$ in the above trivialization.  Write
\[
DF(x) = \begin{pmatrix}
1+\phi'(x) & \lambda \\
\phi'(x)   & \lambda
\end{pmatrix}.
\]
For any $(u,v)^{\mathsf{T}}$,
\[
DF(x)\begin{pmatrix}u\\ v\end{pmatrix}
= \begin{pmatrix}
(1+\phi'(x))u + \lambda v \\
\phi'(x)u + \lambda v
\end{pmatrix}
=: \begin{pmatrix} U \\ V \end{pmatrix}.
\]
Using the relations obtained from the invariance equation,
\[
g'(x) = 1+\phi'(x)+\lambda\Psi'(x),\qquad
\Psi'(g(x)) = \frac{\phi'(x)+\lambda\Psi'(x)}{g'(x)},
\]
a direct computation yields
\begin{align*}
\pi_{g(x)}(U,V) &= V - U\,\Psi'(g(x)) \\
&= \bigl(\phi'u+\lambda v\bigr) - \bigl((1+\phi')u+\lambda v\bigr)\Psi'(g(x)) \\
&= \Bigl[\phi' - (1+\phi')\Psi'(g(x))\Bigr]u + \lambda\bigl[1-\Psi'(g(x))\bigr]v \\
&= -\frac{\lambda\Psi'(x)}{g'(x)}\,u + \frac{\lambda}{g'(x)}\,v
 = \frac{\lambda}{g'(x)}\bigl(v-u\Psi'(x)\bigr) \\
&= \frac{\lambda}{g'(x)}\,\pi_x(u,v).
\end{align*}
Therefore $\overline{DF}(x)$ acts as multiplication by $\lambda/g'(x)$ on $N_x\simeq\mathbb{R}$.

\medskip

By definition of $r$-normal hyperbolicity, there exists a $DF$-invariant continuous splitting
\[
T_\Gamma(\mathbb{T}\times\mathbb{R}) = T\Gamma \oplus E^s,
\]
where $E^s$ is one-dimensional and $DF$-invariant with uniform contraction.  For each $x$,
the restriction of the projection $\pi_x$ to the stable fiber $E^s_x$,
\[
\pi_x|_{E^s_x} : E^s_x \longrightarrow N_x,
\]
is a linear isomorphism because $E^s_x \cap T_x\Gamma = \{0\}$ and both spaces are
one-dimensional.  Hence we can select a unique continuous section $w(x)\in E^s_x$ such that
\[
\pi_x(w(x)) = 1 \quad \text{for all } x\in\mathbb{T}.
\]
Write $DF(x)w(x) = \nu(x)\, w(g(x))$, where $\nu(x)$ is the stable multiplier.
Applying $\pi_{g(x)}$ to this identity and using the commutativity
$\pi_{g(x)}\circ DF(x) = \overline{DF}(x)\circ\pi_x$, we obtain
\[
\frac{\lambda}{g'(x)} = \overline{DF}(x)(1)
= \overline{DF}(x)\bigl(\pi_x(w(x))\bigr)
= \pi_{g(x)}\bigl(DF(x)w(x)\bigr)
= \nu(x)\,\pi_{g(x)}(w(g(x)))
= \nu(x).
\]
Thus the stable multiplier is precisely $\nu(x)=\lambda/g'(x)$.

\medskip

The definition of $r$-normal hyperbolicity  implies that there exists a
Riemannian metric on $TM|_\Gamma$ such that for every $x$,
\[
\|DF(x)|_{E^s_x}\| \; \bigl\|(DF(x)|_{T_x\Gamma})^{-1}\bigr\|^{\,r} < 1,
\qquad
\|DF(x)|_{E^s_x}\| < 1.
\]
Because both $T_x\Gamma$ and $E^s_x$ are one-dimensional, the operator norm of a linear
map on a one-dimensional space equals the absolute value of the corresponding multiplier.
Consequently,
\[
|\nu(x)| \cdot |\tau(x)|^{-r} < 1,
\]
i.e.\
\[
\frac{\lambda}{g'(x)} \Bigl(\frac{1}{g'(x)}\Bigr)^{\!r} < 1.
\]
This gives $g'(x)^{r+1} > \lambda$.

		
		\medskip
		\noindent\textbf{Sufficiency ($\Leftarrow$).}
		Now assume $g'(x)>\lambda^{1/(r+1)}$ for every $x$.
		Set $\tau=\lambda^{1/(r+1)}$ and $\theta(x)=\lambda/g'(x)$. Then
		$\theta(x)<\lambda/\tau\le \tau< g'(x)$ and $\theta(x)\,g'(x)^{-r}<1$.
		
		\medskip
	
		Take the basis $e_1(x)=(1,\Psi'(x))^{\mathsf{T}}$, $e_2(x)=(0,1)^{\mathsf{T}}$.
		Direct computation gives
		\[
		DF(x)e_1(x)=g'(x)e_1(g(x)),\qquad
		DF(x)e_2(x)=\lambda e_1(g(x))+\frac{\lambda}{g'(x)}e_2(g(x)).
		\]
		Hence the matrix of $DF(x)$ in $(e_1,e_2)$ is
		\[U(x)=\begin{pmatrix} g'(x) & \lambda \\ 0 & \theta(x) \end{pmatrix}.\]
		Look for a stable vector of the form $w(x)=\tilde a(x)e_1(x)+e_2(x)$ with
		$DF(x)w(x)=\theta(x)w(g(x))$. This yields the functional equation
		\[
		\tilde \alpha(g(x))=\frac{(g'(x))^2}{\lambda}\tilde \alpha(x)+g'(x)
		\;\Longleftrightarrow\;
		\tilde \alpha(x)=\frac{\lambda}{(g'(x))^2}\tilde \alpha(g(x))-\frac{\lambda}{g'(x)}.
		\]
		Note $\tau=\lambda^{1/(r+1)}$, $\lambda/\tau^2=\lambda^{1-2/(r+1)}\le1$. Because \[\displaystyle\kappa:=\sup_x\frac{\lambda}{(g'(x))^2}
		<\frac{\lambda}{\tau^2}\le1 \quad \text{for all } r\ge 1, \]
		the operator
		\[(\mathcal{Q}\tilde{\alpha})(x)=\frac{\lambda}{(g'(x))^2}\tilde{\alpha}(g(x))-\frac{\lambda}{g'(x)}\]
		is a strict contraction on $C^0(\mathbb{T})$.
		By the Banach fixed-point theorem there exists a unique continuous $\tilde \alpha$
		solving the equation.
		Define $E^s_x=\operatorname{span}\{w(x)\}$ with $w(x)=\tilde \alpha(x)e_1(x)+e_2(x)$.
		Then $TM|_\Gamma=T\Gamma\oplus E^s$ is a $DF$-invariant splitting and
		$DF|_{E^s_x}$ is multiplication by $\theta(x)$.
		
		\medskip

		The vector fields $e_1(x)$ and $w(x)$ form a continuous frame of the tangent bundle
$TM|_\Gamma$.  We define a continuous
Riemannian metric by requiring $\{e_1(x), w(x)\}$ to be orthonormal.  Concretely, for
any $u,v\in T_{(x,\Psi(x))}M$ written uniquely as
$u = a_1 e_1(x) + a_2 w(x)$,\; $v = b_1 e_1(x) + b_2 w(x)$, set
\[
\langle u,v\rangle_x = a_1 b_1 + a_2 b_2.
\]
This form is symmetric, positive definite, and varies continuously with $x$ because the
frame depends continuously on $x$ and the coordinate expressions are continuous linear
isomorphisms on each fibre.  In this metric,
\begin{align*}
&\|DF|_{T_x\Gamma}\| = \|DF(x) e_1(x)\| = \|g'(x) e_1(g(x))\| = g'(x),\\
&\|DF|_{E^s_x}\| = \|DF(x) w(x)\| = \|\theta(x) w(g(x))\| = \theta(x),\\
&\|(DF|_{T_x\Gamma})^{-1}\| = \bigl\| \tfrac{1}{g'(x)} e_1(x) \bigr\| = \frac{1}{g'(x)}.
\end{align*}
		Thus
		\[
		\|DF|_{E^s_x}\|\;\bigl\|(DF|_{T_x\Gamma})^{-1}\bigr\|^r
		= \frac{\lambda}{(g'(x))^{r+1}} < 1,\qquad
		\|DF|_{E^s_x}\| = \frac{\lambda}{g'(x)} < 1,
		\]
		uniformly in $x$. This means $\Gamma$ is $r$-normally hyperbolic.
	\end{proof}

\end{document}
